\begin{figure}[htbp]
\centering
\begin{tikzpicture}
\begin{axis}[
    width=12cm, height=8cm,
    xlabel={$x$},
    ylabel={$f(x)$},
    xmin=-0.5, xmax=5.5,
    ymin=-0.3, ymax=3.5,
    grid=major,
    grid style={UofSC50Black!30},
    axis lines=middle,
    xtick={1,2,3,4,5},
    ytick={1,2,3},
    clip=false,
    legend pos=north east
]

% Define the function: f(x) = 0.5 + 0.8*sin(0.8*x) + 0.3*x
% Integration bounds
\pgfmathsetmacro{\aval}{1}
\pgfmathsetmacro{\bval}{4.5}
\pgfmathsetmacro{\nrect}{7}
\pgfmathsetmacro{\dx}{(\bval-\aval)/\nrect}

% Riemann sum rectangles (left endpoints)
\foreach \i in {0,1,...,6} {
    \pgfmathsetmacro{\xleft}{\aval + \i*\dx}
    \pgfmathsetmacro{\xright}{\xleft + \dx}
    \pgfmathsetmacro{\fleft}{0.5 + 0.8*sin(0.8*\xleft r) + 0.3*\xleft}
    \addplot[fill=UofSCAtlantic!25, draw=UofSCAtlantic, thick, forget plot]
        coordinates {(\xleft,0) (\xleft,\fleft) (\xright,\fleft) (\xright,0) (\xleft,0)};
}

% Shaded area under curve (more accurate representation)
\addplot[fill=UofSCGarnet!15, draw=none, domain=\aval:\bval, samples=100, forget plot]
    {0.5 + 0.8*sin(0.8*x r) + 0.3*x} \closedcycle;

% The function curve
\addplot[UofSCBlack, very thick, domain=0:5.2, samples=150] {0.5 + 0.8*sin(0.8*x r) + 0.3*x};
\addlegendentry{$f(x)$}

% Markers for bounds a and b
\node[below, font=\bfseries, UofSCGarnet] at (axis cs:\aval,-0.15) {$a$};
\node[below, font=\bfseries, UofSCGarnet] at (axis cs:\bval,-0.15) {$b$};

% Vertical lines at bounds
\draw[UofSCGarnet, thick, dashed] (axis cs:\aval,0) -- (axis cs:\aval,{0.5 + 0.8*sin(0.8*\aval r) + 0.3*\aval});
\draw[UofSCGarnet, thick, dashed] (axis cs:\bval,0) -- (axis cs:\bval,{0.5 + 0.8*sin(0.8*\bval r) + 0.3*\bval});

% Delta x label
\draw[UofSCHorseshoe, thick, |<->|] (axis cs:{1+\dx},-0.22) -- (axis cs:{1+2*\dx},-0.22);
\node[below, UofSCHorseshoe, font=\small\bfseries] at (axis cs:{1+1.5*\dx},-0.22) {$\Delta x$};

% Rectangle height label
\pgfmathsetmacro{\xmid}{1 + 2*\dx}
\pgfmathsetmacro{\fmid}{0.5 + 0.8*sin(0.8*\xmid r) + 0.3*\xmid}
\draw[UofSCCongaree, thick, <->] (axis cs:{\xmid+0.05},0.05) -- (axis cs:{\xmid+0.05},{\fmid-0.05});
\node[right, UofSCCongaree, font=\small\bfseries] at (axis cs:{\xmid+0.1},{\fmid/2}) {$f(x_i)$};

\end{axis}

% Integral notation annotation
\node[anchor=north west, align=left, font=\small] at (7.5,6.5) {
    \textcolor{UofSCGarnet}{\textbf{Definite Integral:}}\\[4pt]
    $\displaystyle\int_a^b f(x)\,\dd x = \lim_{n \to \infty} \sum_{i=1}^{n} f(x_i)\,\Delta x$\\[12pt]
    \textcolor{UofSCAtlantic}{\textbf{Riemann Sum:}}\\[4pt]
    $\displaystyle\sum_{i=1}^{n} f(x_i)\,\Delta x \approx \int_a^b f(x)\,\dd x$\\[8pt]
    where $\Delta x = \dfrac{b-a}{n}$
};

\end{tikzpicture}
\caption{Visualization of the definite integral as area under a curve. The function $f(x)$ is shown in black. The shaded Garnet region represents the exact area $\int_a^b f(x)\,\dd x$. The Atlantic blue rectangles show a Riemann sum approximation with $n=7$ rectangles of width $\Delta x = (b-a)/n$. As $n \to \infty$, the Riemann sum converges to the definite integral.}
\label{fig:integral_area}
\end{figure}
